\section{Introduction} \label{sec:introduction}
Foams are among the most familiar yet physically complex systems encountered in everyday life. From the froth on your morning coffee to the structure of biological tissues, foams occupy a unique place in the study of soft matter. At their core, a foam is a dispersion of bubbles seperated by a liquid medium, and the collective behaviour of these bubbles; how they coarsen over time and how the foam restructures because of this, is they core of what is being investigated in this report.\\

Coarsening is a fundamental property of these foams. It is the process by which bubbles exchange gas across their shared interfaces causing larger bubbles to get larger and smaller bubbles get smaller. Over time, the average bubble area increases and the number of bubbles decrease. During this coarsening, bubbles gain new neighbours as the smaller ones disappear causing topological rearrangements in the foam. Understanding the statistical character of the individual bubble trajectories, and how random they are, within this evolving foam structure is the main motivation of this project.\\

In this project, we use data from a 2D foam simulated by Prof. Simon Cox. In this data we are given the positions of the bubble centres at each timestep, along with individual properties such as area and coordination number. Rather than characterising the system through bulk statistical properties, we apply tools borrowed from econophysics to investigate the dynamics of this foam. The central analogy is between the bubble's displacements through the foam and the price movements in a financial market. In both cases, one is interested in the nature of the randomness involved. Whether it is completely random with simple Brownian diffusion, and whether there is a need to characterise a different type of randomness.\\

Brownian motion is used in this project as a baseline of randomness for comparison. It is well understood in it's statistical properties, Gaussian displacement distributions, mean squared displacement that scales linearly in time, and vanishing temporal autocorrelation.This means we can use it as a sanity check for all of our methods, if we apply our methods to it and get the expected results, we can be more confident in our results when we apply the same methods to the bubble trajectories. Deviations from this baseline are of particular interest as they suggest the presence of underlying statistical properties of the foam that influence individual bubble motion.\\

The report is structured as follows: Section \ref{sec:theoreticalbackground} reviews the relevant foam physics, the statistical mechanics of Brownian motion, and the tools from econophysics we will be using; log-returns, volatility, temporal autocorrelation, etc. Section \ref{sec:compmethod} describes the use of pandas to organise the data in a data-frame, extract information, and how we apply computational methods to this information. In section \ref{sec:resultsanddiscussion} we present our results and discuss their implications for understanding bubble motion in coarsening foams. Finally, in section \ref{sec:conclusion} we summarise our findings and suggest directions for future research.
